\documentclass[10pt,a4paper]{amsart}
\usepackage[margin=19mm]{geometry}
\usepackage{amsmath,amssymb,mathtools}
\usepackage[scr=rsfs,cal=euler]{mathalfa}
\usepackage{xfrac}
\usepackage[unicode,hidelinks,bookmarks=false]{hyperref}
\newcommand{\ie}{i.e.\ }
\newcommand{\cf}{cf.\ }
\newcommand{\resp}{respectively\ }
\newcommand{\Subsection}{\subsection}
\providecommand{\nobreakdash}{\nobreak\hbox{-}}
\setlength{\emergencystretch}{2em}
\multlinegap0pt
\usepackage{amssymb}
\usepackage{xcolor}
\newtheorem{lem}{Lemma}
\newcommand{\om}{\omega}
\newcommand{\sm}{\sigma}
\title{Pureness of Certain Crossed Product C*-Algebras}
\author{Dawn Archey, Julian Buck, Javad Mohammadkarimi, N. Christopher Phillips, Apurva Seth}
\date{Source: arXiv:2604.10080v1 (CC BY 4.0)}
\begin{document}
\maketitle

\noindent\emph{Source and licence.} Pureness of Certain Crossed Product C*-Algebras. Version arXiv:2604.10080v1 (CC BY 4.0), distributed by arXiv under the Creative Commons Attribution 4.0 International licence, http://creativecommons.org/licenses/by/4.0/, as recorded in the arXiv licence metadata for this exact version and shown on the abstract page https://arxiv.org/abs/2604.10080v1. Full licence notice: \texttt{licenses/arxiv-2604.10080-CC-BY-4.0.txt}.\par\smallskip

\noindent\textit{Editorial scope.} Counted unit: the article's lem L\_1\_TracesTensorProd (source0.tex lines 2043--2048). The article's complete original proof of it is delivered with it (source0.tex lines 2050--2108, one \texttt{proof} environment of 59 lines). No mathematics is written, repaired or re-derived: the statement and the proof are the article's own lines, in the article's own notation, and no retained line is substituted or re-flowed.  No further source-local context is needed: every cross-reference of every retained line resolves inside the two delivered spans.  Nothing was omitted from any retained span, so the package carries no omission record.  No retained line contains a citation, so the wrapper carries no bibliography and no bibliography entry.  No retained line contains a figure environment, an image-inclusion command or drawing code, so no image is delivered and the package declares no asset dependency.  Editorial formatting only, in addition: an AMS \texttt{amsart} preamble that restates the article's own declarations and macros used by the retained lines, namely the environments \texttt{lem} on separate counters, together with the standard packages those lines need.  The retained environments print in this document as Lemma 1 in order of appearance, whereas the article prints the same items with the numbering of its own full document.  The complete native source of the article as distributed in the arXiv e-print for this exact version is bundled unchanged as \texttt{sources/198244/source0.tex}.
\begin{lem}\label{L_1_TracesTensorProd}
Let $A$ and $B$ be simple unital C*-algebras.
Assume $B$ has a unique tracial state $\sm$.
Then the map $\rho \mapsto \rho \otimes \sm$
is a bijection from $T(A)$ to $T(A \otimes_{\min} B)$.
\end{lem}

\begin{proof}
Let $\tau$ be a tracial state on $A \otimes_{\min} B$.
For fixed $a \in A_{+}$, define a map $\om_{a}$ on
$B$ by $\om_{a} (b) = \tau (a \otimes b)$.
Then
$\om_{a}$ is a bounded linear functional on $B$,
and is positive.
For any $b,c \in B$, we have
%
\begin{align*}
\om_{a} (bc) - \om_{a}(cb) &= \tau (a \otimes bc) -
\tau (a \otimes cb) \\
&= \tau (a \otimes bc - a \otimes cb) \\
&= \tau ( (a \otimes b) (1_{A} \otimes c) -
(1_{A} \otimes c) (a \otimes b) ) = 0
\end{align*}
%
and so $\om_{a}$ has the trace property.
Therefore,
it is a nonnegative multiple of $\sm$, say
$\om_{a} (b) = \rho (a) \sm (b)$.
Writing
$a \in A$ as a linear combination of positive
elements allows the extension of $\rho$ to a map
$A \to {\mathbb{C}}$.
Positivity and the condition
that $\rho (1) = 1$ are immediate.
The trace
property of $\rho$ is easily verified via the
calculation
%
\[
\rho (xy) = \rho(xy) \sm (1_{B}) = \om_{xy}(1_{B})
= \tau (xy \otimes 1_{B}) = \tau (yx \otimes 1_{B})
= \rho (yx) \sm(1_{B}) = \rho (yx).
\]
%
Additivity of $\rho$ on $A_{+}$ follows from the
observation (using the linearity of $\tau$) that
%
\begin{align*}
\rho (a_{1} + a_{2}) \sm (b)
= \om_{a_{1} + a_{2}} (b)
&= \tau ((a_{1} + a_{2}) \otimes b) \\
&= \tau (a_{1} \otimes b) + \tau (a_{2} \otimes b)
= \rho (a_{1}) \sm (b) + \rho (a_{2}) \sm (b)
\end{align*}
%
for any $a_{1}, a_{2} \in A_{+}$, and linearity
of $\rho$ follows.
Hence $\rho$ is a tracial state
on A.
That $\rho \otimes \sm = \tau$ follows from
density of the linear span of the elementary tensors.
This gives surjectivity of the map
$\rho \mapsto \rho \otimes \sm$,
and injectivity follows easily by evaluation on
$a \otimes 1_{B}$.
\end{proof}

\end{document}
